\chapter{Time series, dependence and computational scale}\label{ch:scale}
\section{Three different large-data regimes}
A bank may observe millions of short customer histories, a smaller panel of
branches exposed to common shocks, or one long aggregate balance-sheet
series. These regimes have different asymptotics and different independent
information. Pooling them into one large table does not make them independent.

For a stationary scalar score series with summable autocovariances, the
variance of its average is approximately
\begin{equation}
 \Var(\bar\phi_T)\simeq
 \frac{1}{T}\left\{\gamma_0+2\sum_{k\geq1}\gamma_k\right\},
 \qquad \gamma_k=\Cov(\phi_t,\phi_{t-k}).
 \label{eq:lrv}
\end{equation}
For an AR(1)-like correlation pattern, the variance-based illustrative
effective sample size is $T(1-\rho)/(1+\rho)$.
At $\rho=0.9$, it is about $T/19$. This calculation concerns a stationary
mean; it is not a general causal-information measure for nonstationary
bank data.

\section{A fold is an information boundary}
Customer-level cross-fitting keeps the complete customer's trajectory
outside its nuisance-training fold. Household or connected-business
interference can require a larger unit. A prospective evaluation also uses
training data available before its evaluation origin, with embargoes covering
feature lookback and overlapping outcome windows where needed.

Retrospective estimation under a justified stationary sampling design may
use different folds from a prospective deployment simulation. The product
must distinguish the two. Randomly splitting account-month rows leaks
individual histories. Pure time splitting can still share a common economic
shock between training and test entities. There is no universal split that
repairs every dependence structure.

Multiway-cluster DML splits along cluster dimensions under a specific
sampling framework \citep{cluster2021}. For crossed clustering, a familiar
score covariance construction adds cluster contributions and subtracts their
intersection, schematically
\[
 \widehat V=\widehat V_{\text{entity}}
            +\widehat V_{\text{date}}
            -\widehat V_{\text{entity,date}}.
\]
Scaling and finite-sample corrections depend on the estimator. Few clusters,
unbalanced groups and indefinite finite-sample estimates require disclosure
and suitable alternatives. Increasing the number of accounts cannot replace
missing independent policy dates.

HAC estimators, block bootstraps and cluster bootstraps must match the score
process and dependence assumptions. The report records bandwidth/block
length, cluster counts, sensitivity and any finite-sample correction.
Generic package defaults are not a scientific inference specification.

\section{Changing regimes and transport}
Rates, underwriting policies and customer behavior change. A pooled
coefficient across pre-pandemic, crisis and high-rate periods may have little
decision relevance. Stability diagnostics should examine treatment policy,
covariate support, residual distributions and effect heterogeneity by regime.
Change-point detection is evidence of distributional change, not a unique
explanation for it.

Transporting a conditional effect from source $S=1$ to target $S=0$
may use $\E[\tau(X)\mid S=0]$ if the relevant conditional response is
invariant and target covariates have source support. The target distribution
must be recorded. Unobserved effect modifiers, product versions and changes
in interference can invalidate transport even when feature distributions
overlap. Banking macro and equilibrium responses can change the mechanism
itself \citep{beraja2019,benetton2021}.

\section{Interference and shared constraints}
Customers can share households, firms can share owners, and offers can
compete for branch capacity. In such settings write
$Y_i(a_i,z_i)$, where $z_i=f_i(\bm a_{-i},\mathcal G)$ summarizes
others' assignments through a declared network $\mathcal G$.
The exposure mapping is a substantive restriction: it asserts which
features of others' treatment matter. Individual potential outcomes
indexed only by $a_i$ are inadequate under interference
\citep{hernan2020}.

Direct effects compare $Y_i(1,z)-Y_i(0,z)$ at fixed exposure;
spillover effects compare exposures at fixed own action.
A group policy effect changes both. Household or branch randomization,
or randomized treatment saturation across groups, can create the relevant
variation when operationally appropriate. Independent-unit inference must
then use the randomized or plausibly independent groups.

For observational networks, homophily, endogenous links and unobserved
shared shocks complicate identification. Neither a graph database nor
clustered standard errors eliminates these problems. The initial product
should support declared partial-interference groups and report when
cross-group spillovers remain plausible. General network identification
is a research extension, not a promised default feature.

\section{Compute around scientific boundaries}
The execution engine should materialize a versioned eligible population,
partition by the independent unit, compute fold-specific nuisances, and
aggregate scores with the declared cluster keys. Map tasks may fit independent
folds or bootstrap replicates; reductions retain sufficient statistics and
covariance blocks. Every retry uses the same input snapshot and random seed.
Hyperparameter search occurs inside nuisance-training partitions.

Dense data footprints are already material. At $10^8$ rows and 256
float32 features, the feature array alone is $102.4$ GB in decimal units,
before identifiers, outcomes, fold indicators, copies and model state.
Use columnar projection, predicate pushdown, partition pruning, streamed
score evaluation and sparse representations when justified. Avoid materializing
all lagged columns and all horizons for every exploratory study.

Finite-state filtering costs $O(NTK^2)$ for dense $K$-state transitions,
or less with sparse structure. The forward pass needs $O(K)$ working
state per active sequence; storing all smoothing messages costs $O(TK)$
per sequence unless checkpointing/recomputation is used. State dimension,
not just row count, can dominate. Particle methods trade computation for
approximation variance and degeneracy diagnostics.

Temporal graph search can be more difficult still. A coarse
$p^2L$ candidate count already gives three million directed lag candidates
for $p=500,L=12$, before conditioning-set search. Apply justified temporal
and domain constraints, variable grouping and staged screening. Screening
can discard true parents and must be evaluated as part of the discovery
procedure. Distributed execution does not eliminate the statistical
multiple-testing or combinatorial problem.

\section{A measurable scale qualification}
The benchmark program should use three tiers: $10^5$ rows for correctness,
$10^7$ for multi-team workflows, and $10^8$ for large group studies. These
are future test sizes, not achieved performance claims. Benchmark fixed
estimands and equivalent estimators on single-node and distributed execution.
Measure peak memory, wall time, compute cost, serialization, retry behavior
and numerical agreement. Report nuisance-fitting and data-preparation costs
separately.

Use an explicitly documented reference machine before stating latency
targets. A candidate target is interactive cohort diagnostics within 30
seconds on a precomputed study snapshot and scheduled large-study completion
within one working day. These are design goals contingent on workload and
hardware. They must not be advertised as measured service levels.
Scientific stopping rules take precedence: a fast job with unsupported
actions or invalid uncertainty has not completed an acceptable study.

\section{Dynamic responses: local projections and structural VARs}
A treasury question often concerns a path: how does a funding shock today
change deposits, lending and contribution over the next twelve months?
Write a stable reduced-form VAR as
\begin{equation}
 X_t=\sum_{\ell=1}^{L}A_\ell X_{t-\ell}+u_t,
 \qquad u_t=C\varepsilon_t,\quad
 \E[\varepsilon_t\varepsilon_t']=I.
 \label{eq:var}
\end{equation}
The lag coefficients describe conditional dynamics. The structural impact
matrix $C$ maps economically interpretable shocks to reduced-form
innovations. Estimating $A_\ell$ and $\Sigma_u=CC'$ does not in general
identify $C$. For any orthogonal $Q$, $(CQ)(CQ)'=\Sigma_u$.
An identifying restriction must choose among observationally compatible
shock decompositions or return a set of responses.

For a VAR(1), the response at horizon $h$ to shock $j$ is
$A^h C e_j$. A recursive timing restriction selects a triangular factor
of $\Sigma_u$. Its interpretation is a statement about which variables
can respond contemporaneously. Ordering funding before lending may be
plausible at one frequency and implausible at another. The selected order
must be attached to the impulse-response report as an assumption, not
hidden in a default Cholesky routine.

Local projections estimate each horizon directly, for example
\begin{equation}
 Y_{t+h}=\alpha_h+\beta_h D_t+\Gamma_h'W_t+v_{t+h}^{(h)},
 \label{eq:lp}
\end{equation}
where $D_t$ is the shock or endogenous intervention variable and $W_t$
contains the declared controls. Without an identification argument,
$\beta_h$ is a projection coefficient. When $D_t$ represents a suitably
identified shock, it can estimate a dynamic causal response. Direct horizon
estimation permits flexible responses and avoids imposing all finite-order
VAR restrictions, often with a variance cost. The review by
\citet{jorda2024} explains both the connection and the finite-sample
tradeoffs; it does not justify a universal preference for one method.

\subsection{External instruments identify a particular shock}
An external instrument $Z_t$ must correlate with the shock of interest and
be orthogonal to the other shocks that enter the dynamic outcome.
Under the moving-average representation
$X_t=\sum_{h\geq0}\Theta_h\varepsilon_{t-h}$, assume
$\E[Z_t\varepsilon_{1t}]=\alpha\ne0$ and
$\E[Z_t\varepsilon_{j,t+k}]=0$ for other contemporaneous shocks
and for the relevant nonzero leads and lags. Normalize the impact of shock
one on $X_{1t}$ to unity. Then
\begin{equation}
 \Theta_{h,i1}=\frac{\E[Z_t X_{i,t+h}]}{\E[Z_t X_{1t}]}.
 \label{eq:lpiv}
\end{equation}
The numerator equals $\alpha\Theta_{h,i1}$ because all other shock
covariances vanish; the denominator equals $\alpha$ under the normalization.
This is the dynamic covariance-ratio argument in
\citet[Section 2.3]{stock2018}. Controls require the corresponding
partialled-out restrictions.

The lead-lag exclusion is substantive. A monetary-policy announcement
surprise might also reveal information about future economic conditions.
A bank funding instrument may reflect a local demand shock. A strong first
stage does not exclude either channel. Direct LP-IV and a proxy SVAR have
different information requirements: the latter normally needs a suitable
invertibility or recoverability argument, while the direct approach in the
stated source can avoid invertibility under stronger lead-lag exclusions.
The product must retain that distinction when offering both estimators.

An intervention that raises a bank's posted rate for three months is also
different from a one-time structural shock with endogenous future rate
responses. An impulse response describes the latter under a maintained
system. A sustained policy path requires a different intervention or
structural policy simulation. Calling both a ``25-basis-point scenario''
would conceal an economically material difference.

\section{Inference concerns the score process}
For a scalar moment estimator satisfying
$T^{-1}\sum_t\psi_t(\widehat\theta)=0$, a first-order expansion gives
\[
 \sqrt T(\widehat\theta-\theta_0)
 \simeq -J^{-1}T^{-1/2}\sum_t\psi_t(\theta_0),
 \qquad J=\E[\partial_\theta\psi_t(\theta_0)].
\]
Under an appropriate dependent-data central limit theorem, its covariance
uses the long-run variance of $\psi_t$, not simply that of outcomes or
regression residuals. With matrix moments the same expression uses
$J^{-1}\Omega J^{-\prime}$. A kernel HAC estimate can be written
\[
 \widehat\Omega=\widehat\Gamma_0+
 \sum_{k=1}^{b}K(k/b)(\widehat\Gamma_k+\widehat\Gamma_k'),
 \quad \widehat\Gamma_k=T^{-1}\sum_{t=k+1}^{T}
     \widehat\psi_t\widehat\psi_{t-k}'.
\]
Bandwidth $b$, kernel, centering, finite-sample adjustment and the range of
available $T$ determine the approximation. Short series and persistent
scores can yield poor finite-sample coverage even with a familiar HAC label.

Overlapping horizons make the distinction useful. In an AR(1), repeated
substitution gives
$X_{t+h}=\rho^h X_t+\sum_{j=1}^{h}\rho^{h-j}u_{t+j}$.
The projection residuals overlap and are serially correlated. Some
lag-augmented local-projection constructions make the relevant score
serially uncorrelated under stronger innovation conditions, allowing
heteroskedasticity-robust inference. That conclusion comes from the
regressor--innovation structure, not from denying residual overlap.
\citet[Section 6]{jorda2024} distinguishes these cases. Its reviewed
uniform-inference result should not be transferred without proof to arbitrary
multiway bank panels or learned nonlinear nuisance scores.

Pointwise intervals for twelve horizons also do not give a simultaneous
95\% band for the response path. A maximum-statistic bootstrap, a justified
joint Gaussian approximation or a conservative multiplicity adjustment
can answer the pathwise question. The resampling must preserve the time
and cluster structure. A funding decision depending on the discounted sum
of the response needs the covariance across horizons, which can be much
more informative than reading individual significance stars.

\section{A repeated-sample demonstration of common shocks}
The supplied experiment draws a score panel
\[
 \phi_{it}=a_i+f_t+\epsilon_{it},\qquad
 a_i,f_t,\epsilon_{it}\ \hbox{independent }N(0,1),
\]
with 96 entities and 36 dates. The target mean is zero. Direct variance
calculation yields
\begin{equation}
 \Var(\bar\phi)=\frac1{96}+\frac1{36}+\frac1{96\cdot36}.
 \label{eq:crossed-variance}
\end{equation}
The resulting true standard error is approximately $0.1962$.
There are 3,456 rows, but shared entity and date shocks dominate the
uncertainty. Treating rows as independent counts those shocks repeatedly.

\begin{table}[htbp]\centering\small
\begin{tabular}{@{}lrr@{}}
\toprule
Variance method & Mean SE & 95\% coverage\\
\midrule
Rows independent & 0.029 & 0.227\\
Entity clusters & 0.103 & 0.693\\
Entity and date clusters & 0.193 & 0.944\\
\bottomrule
\end{tabular}

\caption{Coverage over 1,200 independently generated panels. Critical values
are normal; the crossed construction uses entity and date corrections. This
is a stated Gaussian score design, not a theorem for arbitrary bank panels.}
\label{tab:dependence-coverage}
\end{table}

Entity clustering protects against repeated observations of a customer but
misses the date component in this design. The crossed construction is much
closer to the intended coverage, while only 36 date clusters still make
finite-sample approximation relevant. The code retains the Monte Carlo
standard error of every reported coverage estimate. Adding a million
customers would reduce the first and third terms in
Equation~\eqref{eq:crossed-variance}; it would leave the date term.
This is why the execution planner must preserve the scientific unit rather
than simply distribute all rows uniformly over workers.

\begin{figure}[htbp]\centering
\includegraphics[width=\textwidth]{derivations/dependence_and_irf.pdf}
\caption{Left: dependence changes uncertainty in the repeated score experiment.
Right: two structural timing assumptions produce different responses while
preserving the same reduced-form Gaussian VAR law.}
\label{fig:dependence-irf}
\end{figure}

\section{Large scale has several different bottlenecks}
The scientific workload can be decomposed as
\[
 T_{\rm job}=T_{\rm snapshot}+T_{\rm design}+T_{\rm nuisance}
            +T_{\rm score}+T_{\rm inference}+T_{\rm export},
\]
with separate memory and storage accounts. A fast score reduction can be
irrelevant when source reconstruction or repeated nuisance fitting dominates.
Likewise, a GPU can accelerate dense model fitting while increasing the
cost of copying sparse event data into device memory. The product should
measure each stage and retain the workload configuration with the result.

Consider $N$ independent trajectories, horizon $T$, $K$ latent states,
$p$ observed variables, $L$ candidate lags, $F$ folds and $B$ bootstrap
replicates. These dimensions describe different jobs. Dense HMM smoothing
requires order $NTK^2$ arithmetic per likelihood pass. Refitting all
parameters in every bootstrap can multiply that work by $B$, not merely
add a cheap variance reduction. A temporal graph learner may have
$p^2L$ candidate edges and far more conditioning tests. A flexible
longitudinal nuisance pipeline may fit a distinct learner for each date,
fold and target policy. Reporting only total rows misses each of these
sources of difficulty.

Several computational reductions have scientific conditions. Homogeneous
transition counts can be summed across customers only when the same
parameterization applies. State compression must preserve the decision-
relevant conditional law. Pooling across jurisdictions can improve precision
only under a defensible sharing or hierarchical model. Sparse graph
screening reduces search only by making a restriction on candidate parents.
Each reduction belongs in the model specification and sensitivity study.

The recommended architecture caches reusable source snapshots and
fold-consistent nuisance predictions, streams score evaluation, and reduces
cluster contributions without materializing unnecessary dense copies.
It parallelizes independent fits or sequences before splitting a dependent
sequence. For long-sequence smoothing, checkpointing trades recomputation
for storage; truncating a sequence instead changes the conditional
information available to the estimator. Those are different operations.

\section{A scale result must preserve the statistical result}
The construction chapter supplies a measured streaming reduction workload
and a temporal-revision implementation. These bounded measurements establish
the operation and resource use of those components. They do not establish
end-to-end performance for a causal forest, temporal graph search or
controlled-state likelihood on a banking group's records.

A future large-scale qualification uses identical source snapshots, folds,
target weights and random seeds for a single-node reference and a distributed
implementation. It compares estimates, influence contributions, covariance,
unsupported-cell flags and financial totals, not just runtime. Differences
caused by floating-point summation are bounded; differences caused by changed
folds or dropped records are scientific defects. Retry tests must show that
a failed worker does not duplicate a customer's contribution.

Resource budgets should be expressed per workload: bytes scanned per study,
peak worker memory, total CPU/GPU time, shuffled bytes, model checkpoints,
and the number of complete refits. A claim such as ``supports one hundred
million rows'' becomes useful only after specifying these quantities, the
estimator and the attainable uncertainty. Computation can make a valid
analysis feasible; it cannot create additional independent policy shocks.
